% ---------------------------------------------------------------------------
% Preamble: packages, theorem environments and shared macros.
% Everything that all chapters may rely on lives here.
% ---------------------------------------------------------------------------

\usepackage[T1]{fontenc}
\usepackage[utf8]{inputenc}
\usepackage{lmodern}
\usepackage[english]{babel}
\usepackage[a4paper,margin=3cm]{geometry}

\usepackage{amsmath}
\usepackage{amssymb}
\usepackage{amsthm}
\usepackage{mathtools}

\usepackage{graphicx}
\graphicspath{{figures/}}
\usepackage{xcolor}
\usepackage{tikz}
\usetikzlibrary{arrows.meta,positioning,calc,patterns}

\usepackage{booktabs}
\usepackage{enumitem}
\usepackage{microtype}
\usepackage{algorithm}
\usepackage{algpseudocode}
\usepackage[colorlinks=true,linkcolor=blue!50!black,citecolor=blue!50!black,
            urlcolor=blue!50!black]{hyperref}
\usepackage{cleveref}
\usepackage[numbers,sort&compress]{natbib}

% Bookmarks are plain text: tell hyperref how to spell our macros there.
\pdfstringdefDisableCommands{%
  \def\ER{exists-R}%
  \def\ETR{ETR}%
  \def\Ostar{O*}%
  \def\R{R}%
  \def\N{N}%
  \def\Z{Z}%
  \def\Q{Q}%
  \def\tw{tw}%
  \def\OPT{OPT}%
  \def\lang#1{#1}%
}

% --- Theorem-like environments ---------------------------------------------
\theoremstyle{plain}
\newtheorem{theorem}{Theorem}[chapter]
\newtheorem{lemma}[theorem]{Lemma}
\newtheorem{proposition}[theorem]{Proposition}
\newtheorem{corollary}[theorem]{Corollary}
\newtheorem{claim}[theorem]{Claim}
\newtheorem{conjecture}[theorem]{Conjecture}

\theoremstyle{definition}
\newtheorem{definition}[theorem]{Definition}
\newtheorem{example}[theorem]{Example}
\newtheorem{problem}[theorem]{Problem}
\newtheorem{exercise}[theorem]{Exercise}
\newtheorem{observation}[theorem]{Observation}

\theoremstyle{remark}
\newtheorem{remark}[theorem]{Remark}
\newtheorem*{intuition}{Intuition}

% --- Problem definitions ----------------------------------------------------
% Usage:
%   \begin{probdef}{Vertex Cover}
%     \instance A graph $G=(V,E)$ and an integer $k$.
%     \question Is there a set $S \subseteq V$ with $|S| \le k$ ...
%   \end{probdef}
\newenvironment{probdef}[1]
  {\begin{quote}\textsc{#1}\par\nopagebreak\begin{description}[leftmargin=2em,
     style=nextline,font=\normalfont\itshape]}
  {\end{description}\end{quote}}
\newcommand{\instance}{\item[Instance:]}
\newcommand{\question}{\item[Question:]}
\newcommand{\parameter}{\item[Parameter:]}
\newcommand{\goal}{\item[Goal:]}

% --- Shared notation --------------------------------------------------------
\newcommand{\N}{\mathbb{N}}
\newcommand{\Z}{\mathbb{Z}}
\newcommand{\Q}{\mathbb{Q}}
\newcommand{\R}{\mathbb{R}}
\newcommand{\ER}{\exists\R}                 % the class exists-R
\newcommand{\ETR}{\textsc{ETR}}
\newcommand{\Ostar}{O^{*}}                  % O* notation, suppresses poly factors
\newcommand{\poly}{\mathrm{poly}}
\newcommand{\OPT}{\mathrm{OPT}}
\newcommand{\tw}{\mathrm{tw}}
\newcommand{\lang}[1]{\textsc{#1}}          % problem names in running text
\newcommand{\bigoh}[1]{O\!\left(#1\right)}
\DeclareMathOperator{\dg}{deg}
\DeclareMathOperator{\cost}{cost}
\DeclareMathOperator*{\argmin}{arg\,min}
\DeclareMathOperator*{\argmax}{arg\,max}

% Frequently used problem names.
\newcommand{\Men}{\mathcal{M}}
\newcommand{\Women}{\mathcal{W}}
\newcommand{\pr}[1]{\succ_{#1}}   % w \pr{m} w' reads "m prefers w to w'"
\newcommand{\GS}{\lang{Gale--Shapley}}
\newcommand{\ThreeCol}{\lang{3-Coloring}}
\newcommand{\PlanarThreeCol}{\lang{Planar-3-Coloring}}
\newcommand{\MIS}{\lang{Max-Independent-Set}}

% --- Authoring helpers ------------------------------------------------------
% \todo{...} marks work in progress and is deliberately visible in the PDF.
% \long so that a TODO may span several paragraphs.
\long\def\todo#1{{\color{red}\bfseries[TODO: #1]}}

% A figure that has not been drawn yet: prints a framed grey box describing
% the intended picture, so the book compiles and the gap is visible.
% #1 = height, #2 = description.  See book/figures/README.md.
\newcommand{\figplaceholder}[2]{%
  \setlength{\fboxsep}{6pt}%
  \fcolorbox{gray!60}{gray!8}{%
    \parbox[c][#1][c]{0.86\textwidth}{%
      \centering\small\itshape\color{black!65}%
      [\,figure placeholder\,]\par\vspace{0.6em}#2}}}

% Marks a section that goes beyond the lectures.  The heading carries
% "(non-examinable)" as well; this prints the reason underneath it.
\newcommand{\nonexam}[1]{%
  {\small\itshape\noindent\textbf{Not examinable.}\ #1\par}\smallskip}

% Every chapter closes with the list of open-pool questions it is responsible
% for.  Usage:
%   \begin{poolquestions}
%     \poolq{1}{Stable Matchings} Covered in \cref{sec:sm-gs}.
%   \end{poolquestions}
\newenvironment{poolquestions}
  {\section{Open pool questions covered by this chapter}%
   \begin{itemize}[leftmargin=*]}
  {\end{itemize}}
\newcommand{\poolq}[2]{\item \textbf{Pool question #1 (#2).}\ }
